\section{4.4.Normalizer}\label{normalizer}

\subsection{4.4. Módulo de
normalización}\label{muxf3dulo-de-normalizaciuxf3n}

Uno de los principales problemas identificados durante el estudio de las
herramientas existentes fue la heterogeneidad de los formatos utilizados
para representar las vulnerabilidades detectadas.

Slither y Mythril generan estructuras JSON completamente diferentes,
tanto en la forma de representar la severidad como en la descripción de
los hallazgos o la identificación de las localizaciones afectadas.

Esta situación dificulta la comparación directa de resultados y hace
necesaria una fase previa de transformación antes de poder aplicar
mecanismos de correlación.

Para resolver este problema se desarrolló un módulo de normalización
encargado de transformar las salidas heterogéneas en una estructura
común.

\subsubsection{4.4.1. Objetivos del proceso de
normalización}\label{objetivos-del-proceso-de-normalizaciuxf3n}

El proceso de normalización persigue varios objetivos:

\begin{itemize}
\tightlist
\item
  unificar la representación de vulnerabilidades;
\item
  facilitar la correlación entre herramientas;
\item
  preservar las evidencias originales;
\item
  simplificar la generación posterior de informes;
\item
  desacoplar el resto de módulos de las particularidades de cada
  herramienta.
\end{itemize}

Gracias a ello, los componentes posteriores pueden operar sobre una
estructura homogénea independientemente del origen de los resultados.

\subsubsection{4.4.2. Estructura común de los
hallazgos}\label{estructura-comuxfan-de-los-hallazgos}

Cada vulnerabilidad normalizada se representa mediante un conjunto común
de atributos.

Entre los campos principales se encuentran:

\begin{itemize}
\tightlist
\item
  título;
\item
  descripción;
\item
  severidad;
\item
  categoría;
\item
  contrato afectado;
\item
  función afectada;
\item
  localización;
\item
  identificador SWC;
\item
  herramienta de origen;
\item
  evidencia original.
\end{itemize}

La inclusión de la salida original permite mantener la trazabilidad del
análisis y facilita la revisión manual de los resultados.

!TODO: insertar diagrama UML del objeto Finding.

!TODO: revisar nombres exactos de los atributos.

\subsubsection{4.4.3. Normalización de resultados de
Slither}\label{normalizaciuxf3n-de-resultados-de-slither}

La función de normalización correspondiente a Slither transforma la
estructura JSON generada por la herramienta y extrae la información
relevante necesaria para construir los hallazgos comunes utilizados por
EVMAudit.

Durante esta fase se realiza además la asociación de los detectores
propios de Slither con los identificadores SWC correspondientes.

Este proceso permite traducir la terminología específica de Slither a
una representación independiente de la herramienta utilizada.

!TODO: insertar fragmento de código de
\texttt{normalize\_slither\_output()}.

\subsubsection{4.4.4. Normalización de resultados de
Mythril}\label{normalizaciuxf3n-de-resultados-de-mythril}

De forma análoga, la salida generada por Mythril es procesada para
obtener una representación equivalente a la utilizada por Slither.

La transformación permite unificar:

\begin{itemize}
\tightlist
\item
  niveles de severidad;
\item
  nombres de vulnerabilidades;
\item
  identificadores SWC;
\item
  información contextual del hallazgo.
\end{itemize}

Gracias a ello, ambas herramientas pueden ser tratadas posteriormente de
forma transparente por el módulo de correlación.

!TODO: insertar fragmento de código de
\texttt{normalize\_mythril\_output()}.

\subsubsection{4.4.5. Clasificación de
vulnerabilidades}\label{clasificaciuxf3n-de-vulnerabilidades}

Además de unificar la representación de los hallazgos, el módulo de
normalización clasifica las vulnerabilidades detectadas siguiendo las
categorías definidas durante el estudio del estado del arte.

Entre las categorías utilizadas se encuentran:

\begin{itemize}
\tightlist
\item
  vulnerabilidades de control de acceso;
\item
  vulnerabilidades económicas;
\item
  vulnerabilidades relacionadas con la lógica de negocio;
\item
  vulnerabilidades asociadas a la ejecución del contrato.
\end{itemize}

Esta clasificación permite mantener la coherencia entre la parte teórica
del trabajo y la implementación desarrollada.

!TODO: referencia cruzada a la Sección 2.X (clasificación de
vulnerabilidades).

\subsubsection{4.4.6. Preservación de
evidencias}\label{preservaciuxf3n-de-evidencias}

Uno de los principios adoptados durante el diseño de EVMAudit fue evitar
la pérdida de información durante las fases intermedias del análisis.

Por este motivo, cada hallazgo normalizado conserva una referencia a la
salida original generada por la herramienta correspondiente.

Esta decisión facilita:

\begin{itemize}
\tightlist
\item
  la revisión manual por parte del auditor;
\item
  la trazabilidad del proceso;
\item
  la depuración de errores;
\item
  la reproducibilidad de los resultados.
\end{itemize}

La conservación de evidencias resulta especialmente importante en
entornos de auditoría y análisis forense, donde es necesario justificar
el origen de cada vulnerabilidad identificada.

!TODO: insertar diagrama del proceso de normalización.

!TODO: insertar referencia a la Figura X.
